\section{Locking Strategies} 

\subsection{Fine-Grained Locking}

\begin{remark}
    Instead of locking everything coarse-grained (the whole data structure) we can lock individual parts of it. This is called \textbf{fine-grained locking} and allows for higher concurrency.
\end{remark}

\subsection{Read-Write Locks}

\begin{remark}
    With coarse- and fine-grained locking we have complete mutual exclusion on resources but what if we only need exclusive access for writing?
\end{remark}

\begin{definition}
    A \textbf{read-write lock} allow fair access to resources by allowing one writer (every N reads) or multiple readers at the same time.
\end{definition}

\begin{remark}
    To implement a RWL we need to keep track of the number of readers waiting, the number of writers waiting and the number of readers a writer has to wait until it can write (this ensures fairness).
\end{remark}

\subsection{Optimistic Locking}

\begin{remark}
    When for example traversing a list, finding a node with \textbf{hand-over-hand locking} requires locking all nodes on the way which consumes time and energy. But how high is it that two threads edit the same node at the same time? 
\end{remark}

\begin{definition}
    \textbf{Optimistic Locking: } Instead of locking every node we could traverse the list without locking and only lock the node we are interested in. If we lock and notice that the node has been modified and our locks on (predecessor/successor) are nolonger valid we restart the operation.
\end{definition}

\begin{remark}
    Optimistic locking allows us for less overhead on locking but might require multiple traversals if there is high contention.
\end{remark}

\subsection{Lazy Locking}

\begin{definition}
    \textbf{Lazy locking} is a technique where data-elements are not immediately updated, but rather marked as deleted (or inserted) and only updated when a thread traverses the data structure again. This saves the overhead of multiple traversals.
\end{definition}

\subsection{Lock-Free Data Structures}

\begin{remark}
    Using atomic operations (mostly CAS) we can implement data structures without locks. 
\end{remark}

\begin{codeexample}
class ConcurrentCounter {
    private AtomicLong value;
    
    public ConcurrentCounter() {
        value = new AtomicLong(0);
    }
    
    public void increment() {
        long expected;
        do {
            expected = value.get();
        } while (!value.compareAndSet(expected, expected + 1));
    }
    
    public long getValue() {
        return value.get();
    }
}
\end{codeexample}

\begin{important}
    The \textbf{ABA problem} is a common issue in lock-free programming where an atomic operation (like compare-and-swap) succeeds because the value has changed twice (e.g., A -> B -> A), making it appear unchanged to the operation. This can lead to incorrect program states.
\end{important}
